% Moderne Fußnoten-Konfiguration für LuaLaTeX
% Optimiert für bessere Lesbarkeit und moderne Typografie

% --- Footmisc-Paket mit erweiterten Optionen ---
\usepackage[
  bottom,                      % Fußnoten am Seitenende (Standard)
  hang,                        % Hängende Einrückung
  stable,                      % Stabile Referenzen
  multiple,                    % Mehrere Fußnoten pro Stelle möglich
  norule                       % Moderne Option: keine Linie (optional)
]{footmisc}

% --- Grundlegende Abstände ---
\setlength{\footnotemargin}{12pt}        % Etwas mehr Platz (war 10pt)
\setlength{\footnotesep}{8pt}            % Abstand zwischen Fußnoten
\setlength{\skip\footins}{12pt plus 4pt minus 2pt}  % Flexibler Abstand zum Text

% --- Mehrzeilige Fußnoten (aktiviert) ---
\renewcommand{\hangfootparindent}{1.5em}  % Einrückung 2. Absatz
\renewcommand{\hangfootparskip}{6pt}      % Abstand zwischen Absätzen

% --- Erweiterte Formatierung ---
% Fußnotennummer-Stil
\renewcommand{\thefootnote}{\textcolor{secondary}{\arabic{footnote}}}

% Layout der Fußnotennummer im Text (Superscript) - vereinfacht
% \renewcommand{\@makefnmark}{%
%   \hbox{\@textsuperscript{\normalfont\textcolor{primary}{\@thefnmark}}}%
% }

% Layout der Fußnotennummer in der Fußnote selbst - vereinfacht  
% \renewcommand{\@makefntext}[1]{%
%   \parindent 1em
%   \noindent
%   \hb@xt@1.8em{\hss\textcolor{secondary}{\@thefnmark}\,}#1%
% }

% --- Alternative: Moderne Fußnoten mit besserer Trennung ---
% Ersetze Standard-Linie durch moderne Variante
\renewcommand{\footnoterule}{%
  \kern-3pt
  {\color{primary}\hrule width 0.4\columnwidth height 0.5pt}
  \kern 2.6pt
}

% --- Fußnoten-Schriftart (bereits in TocStyle definiert, hier verfeinert) ---
\renewcommand{\footnotesize}{\fontsize{9}{11}\selectfont}  % Etwas größer als Standard
\addtokomafont{footnote}{\color{darkgray}}                  % Dezente Farbe

% --- Erweiterte Fußnoten-Features ---
% Für sehr lange Fußnoten: bessere Umbrüche
\tolerance=5000
\emergencystretch=3pt

% Verhindere Fußnoten-Waisen (einzelne Zeile auf neuer Seite)
\interfootnotelinepenalty=10000

% --- Fußnoten in verschiedenen Umgebungen ---
% Für Tabellen und Minipages (Reihenfolge wichtig!)
\usepackage{footnotebackref}             % Rückverweise (MUSS vor tablefootnote!)
\usepackage{tablefootnote}               % Fußnoten in Tabellen

% --- Spezielle Fußnoten-Stile ---
% Für verschiedene Dokumentbereiche
\newcommand{\symbolfootnote}[1]{%
  % Verwendet Symbole statt Zahlen
  \renewcommand{\thefootnote}{\fnsymbol{footnote}}
  \footnote{#1}
  \renewcommand{\thefootnote}{\arabic{footnote}}
}

% Für Titelseite (oft andere Nummerierung)
\newcommand{\titlefootnote}[1]{%
  \let\thefootnoteold\thefootnote
  \renewcommand{\thefootnote}{\alph{footnote}}
  \footnote{#1}
  \let\thefootnote\thefootnoteold
}

% --- Integration mit Hyperref ---
\makeatletter
\@ifpackageloaded{hyperref}{%
  % Hyperref-kompatible Fußnoten
  \renewcommand{\@makefntext}[1]{%
    \parindent 1em
    \noindent
    \hb@xt@1.8em{\hss\hyper@makecurrent{Hfootnote}%
    \hypertarget{Hfootnote.\theHfootnote}%
    {\textcolor{secondary}{\@thefnmark}}\,}#1%
  }
}{}
\makeatother

% --- Mehrsprachige Fußnoten ---
\makeatletter
\@ifpackageloaded{polyglossia}{%
  % Anpassung für deutsche/englische Texte
  \newcommand{\englishfootnote}[1]{%
    \begin{otherlanguage}{english}
    \footnote{#1}
    \end{otherlanguage}
  }
}{}
\makeatother

% --- Debugging und erweiterte Kontrolle ---
% Zeige Fußnoten-Informationen (für Entwicklung)
\newif\ifdebugfootnotes
% \debugfootnotestrue

\ifdebugfootnotes
\AtBeginDocument{%
  \typeout{=== Fußnoten-Debug ===}
  \typeout{Margin: \the\footnotemargin}
  \typeout{Sep: \the\footnotesep}
  \typeout{Skip footins: \the\skip\footins}
  \typeout{===================}
}
\fi

% --- Praktische Utility-Befehle ---
% Quellenangaben in Fußnoten
\newcommand{\sourcefootnote}[1]{%
  \footnote{\textit{Quelle:} #1}
}

% Anmerkungen des Autors
\newcommand{\authornote}[1]{%
  \footnote{\textit{Anm. d. Verf.:} #1}
}

% Übersetzungshinweise
\newcommand{\transnote}[1]{%
  \footnote{\textit{Übersetzung d. Verf.:} #1}
}

% --- Alternative: Randbemerkungen statt Fußnoten ---
% Für moderne, saubere Layouts
\newcommand{\marginnote}[1]{%
  \marginpar{\footnotesize\raggedright\textcolor{secondary}{#1}}
}

% --- KOMA-Script Integration ---
% Nutzt KOMA's erweiterte Fußnoten-Features
\KOMAoptions{
  footnotes=multiple,          % Mehrere Fußnoten
  footnotes=nomultiple         % oder: keine mehrfachen (je nach Bedarf)
}

% --- Conditional Compilation: Druck vs. Digital ---
\makeatletter
\@ifundefined{ifprintversion}{}{%
  \ifprintversion
    % Druck: Konservativere Fußnoten
    \renewcommand{\footnoterule}{%
      \kern-3pt
      {\color{printgray}\hrule width 0.4\columnwidth height 0.4pt}
      \kern 2.6pt
    }
    \addtokomafont{footnote}{\color{printblack}}
  \else
    % Digital: Moderne, farbige Fußnoten (wie oben definiert)
  \fi
}
\makeatother

% --- Debug-Information ---
\AtBeginDocument{%
  \typeout{=== Fußnoten konfiguriert ===}
  \typeout{Paket: footmisc (erweitert)}
  \typeout{Stil: Hängend mit Farben}
  \typeout{Margin: \the\footnotemargin}
  \typeout{Features: Mehrsprachig, Hyperref-kompatibel}
  \typeout{===============================}
}